\renewcommand{\headline}{Penta$\cdot$खेल - हिंदी}
\renewcommand{\tocent}{हिंदी}
\renewcommand{\tocline}{हिंदी में नियम}
\renewcommand{\translator}{क्लॉड द्वारा प्रारूप, समीक्षा आवश्यक}
\renewcommand{\general}{
एक ऐसा खेल जिसे एक मिनट में समझाया जा सकता है, लेकिन जो बरसों तक दिलचस्प बना रहता है।

दो खिलाड़ियों के बीच यह खेल सिर्फ़ 20 मिनट में ख़त्म हो जाता है। तीन या चार खिलाड़ियों के साथ इसमें 90 मिनट तक लग सकते हैं।

इसमें पासे का इस्तेमाल नहीं होता। यह 5 साल और उससे ऊपर के सभी लोगों के लिए उपयुक्त है।

डिब्बे में हैं: 4$\times$5 रंगीन गोटियाँ, 5 काले और 5 स्लेटी ब्लॉक, और बोर्ड।
}

\renewcommand{\choosext}{अपनी गोटियाँ चुनें}
\renewcommand{\choosex}{
हर खिलाड़ी के पास पांच गोटियाँ होती हैं। डिब्बे में हर खिलाड़ी के लिए गोटियों का एक सेट है।

आपके पास एक नीली, एक लाल, एक सफ़ेद, एक हरी और एक पीली गोटी होती है।

गोटियाँ बोर्ड के उन पांच कोनों से शुरू होती हैं जो उनके रंग से मेल खाते हैं।

वे बीच में मौजूद बड़े चौराहों तक पहुंचना चाहती हैं, जहाँ पहुँचकर वे खेल से बाहर हो जाती हैं।
}

\renewcommand{\setupt}{तैयारी}
\renewcommand{\setup}{
अपनी गोटियों को बाहरी घेरे के बड़े कोनों पर रखें, जो उनके रंग से मेल खाते हों: सफ़ेद गोटी को घेरे के सफ़ेद कोने पर, नीली गोटी को घेरे के नीले कोने पर, और इसी तरह बाकी गोटियों को भी रखें।

काले ब्लॉक बोर्ड के बीचोंबीच स्थित पांचों चौराहों पर रखें।

स्लेटी ब्लॉक बाद में इस्तेमाल करने के लिए बीच में रख दें।
}

\renewcommand{\objectivet}{लक्ष्य}
\renewcommand{\objective}{
सफ़ेद गोटियाँ बीच के सफ़ेद चौराहे तक जाना चाहती हैं, नीली गोटियाँ नीले चौराहे तक, और इसी तरह बाकी भी। मंज़िल हमेशा शुरुआती बिंदु के ठीक सामने, पंचभुज के केंद्र में स्थित बड़ा रंगीन चौराहा होती है।

जीतने के लिए सबसे पहले अपनी \textbf{तीन} गोटियाँ उनकी मंज़िल तक पहुँचाएँ।
}

\renewcommand{\rulest}{नियम}
\renewcommand{\rules}{
अपनी किसी भी गोटी को तारे की रेखाओं या बाहरी घेरे पर, किसी भी दिशा में, जितना चाहें उतना दूर तक चलाएँ।

आप बिना रुके किसी भी खाली कोने या चौराहे पर मुड़ सकते हैं।

आप कभी भी किसी के ऊपर से नहीं फांद सकते\textemdash न ब्लॉक के ऊपर से, न किसी और गोटी के ऊपर से।

\myskip

लेकिन आप किसी ऐसे स्थान पर \textbf{ज़रूर} जा सकते हैं जहाँ पहले से कुछ हो:

\myskip

अगर ऐसा करने से आप किसी काले ब्लॉक को हटाते हैं, तो उसे अपनी पसंद के किसी भी खाली स्थान पर रख दें।

अगर ऐसा करने से आप किसी और गोटी वाले स्थान पर पहुँचते हैं, तो उसके साथ अपनी जगह बदल लें।

\myskip

आप अपनी दो गोटियों की जगह आपस में बदल सकते हैं।

अगर आप ऐसे स्थान पर पहुँचें जहाँ कई गोटियाँ हों, तो उनमें से किसी एक के साथ जगह बदल लें।

\myskip

लगातार दो बार एक जैसी चाल न चलें।

\myskip

जो गोटी अपनी मंज़िल तक पहुँच जाए, उसे बोर्ड से हटाकर बीच में रख दें। इसके बाद स्लेटी ब्लॉकों में से एक उठाएँ और उसे अपनी पसंद के किसी खाली स्थान पर रख दें।

अगर आप इस तरह किसी स्लेटी ब्लॉक को हटाते हैं, तो उसे फिर से बोर्ड से हटा लें।

अगर स्लेटी ब्लॉक ख़त्म हो जाएँ, तो जो पहले से खेल में हैं उनमें से किसी एक को फिर से रखें।

\myskip

अगर किसी दूसरे खिलाड़ी ने आपकी किसी गोटी को उसकी मंज़िल तक पहुँचा दिया है, तो अपनी बारी आते ही आपको उसे हटाना होगा।

\myskip

आख़िरी दौर अंत तक खेला जाता है।
}
